% from ega3-6-fr.tex
\oldpage[III]{163}
de $X'^{(i)}$ par des ouverts affines, \emph{plus fins} respectivement que les
recouvrements $u_i^{-1}(\mathfrak U^{(i)})$ ($i=1,2$). Il est clair que le
bicomplexe
$C^\bullet(\mathfrak U^{(i)},\mathscr P_\bullet^{(i)})
=L_{\bullet\bullet}^{(i)}$ peut être considéré canoniquement comme un
sous-bicomplexe de
$C^\bullet(u_i^{-1}(\mathfrak U^{(i)}),u_i^*(\mathscr P_\bullet^{(i)}))$
($0_{\mathrm I}$, 4.4.3.2); en outre, en choisissant une application
simpliciale (G, II, 5.7) de $\mathfrak U'^{(i)}$ dans
$u_i^{-1}(\mathfrak U^{(i)})$, on définit un homomorphisme de bicomplexes
\[
  C^\bullet(u_i^{-1}(\mathfrak U^{(i)}),u_i^*(\mathscr P_\bullet^{(i)}))
  \longrightarrow
  C^\bullet(\mathfrak U'^{(i)},u_i^*(\mathscr P_\bullet^{(i)}))
\]
d'où, par composition, un homomorphisme de bicomplexes
\[
  L_{\bullet\bullet}^{(i)}
  \longrightarrow
  L_{\bullet\bullet}^{\prime(i)}
  =C^\bullet(\mathfrak U'^{(i)},u_i^*(\mathscr P_\bullet^{(i)})).
\]
En outre, cet homomorphisme est remplacé par un homomorphisme homotope quand
on change d'application simpliciale (G, II, 5.7.1); on a ainsi un
homomorphisme bien défini de foncteurs spectraux :
\[
  {}^{(t)}\mathscr E
  (\mathfrak U^{(1)},\mathfrak U^{(2)};
   \mathscr P_\bullet^{(1)},\mathscr P_\bullet^{(2)})
  \longrightarrow
  {}^{(t)}\mathscr E
  (\mathfrak U'^{(1)},\mathfrak U'^{(2)};
   u_1^*(\mathscr P_\bullet^{(1)}),u_2^*(\mathscr P_\bullet^{(2)})).
  \tag{6.7.10.3}\label{III.6.7.10.3.fr}
\]
On vérifie aussitôt que si $\mathfrak V^{(i)}$ est un recouvrement affine
fini de $X^{(i)}$ plus fin que $\mathfrak U^{(i)}$,
$\mathfrak V'^{(i)}$ un recouvrement affine fini de $X'^{(i)}$, plus fin que
$u_i^{-1}(\mathfrak V^{(i)})$ et que $\mathfrak U'^{(i)}$, le diagramme
\[
\xymatrix@C=34pt@R=30pt{
  C^\bullet(\mathfrak U^{(i)},\mathscr P_\bullet^{(i)})
  \ar[r]\ar[d]
  & C^\bullet(\mathfrak V^{(i)},\mathscr P_\bullet^{(i)})\ar[d]
  \\
  C^\bullet(\mathfrak U'^{(i)},u_i^*(\mathscr P_\bullet^{(i)}))
  \ar[r]
  & C^\bullet(\mathfrak V'^{(i)},u_i^*(\mathscr P_\bullet^{(i)}))
}
\]
est commutatif, ce qui implique que l'homomorphisme (6.7.10.3) ne dépend pas
essentiellement des recouvrements $\mathfrak U^{(i)}$ et
$\mathfrak U'^{(i)}$ considérés. On a donc en fait défini un homomorphisme de
$A$-modules
\[
  {}^{(t)}E
  (X^{(1)},X^{(2)};\mathscr P_\bullet^{(1)},\mathscr P_\bullet^{(2)})
  \longrightarrow
  {}^{(t)}E
  (X'^{(1)},X'^{(2)};
   u_1^*(\mathscr P_\bullet^{(1)}),u_2^*(\mathscr P_\bullet^{(2)})).
  \tag{6.7.10.4}\label{III.6.7.10.4.fr}
\]
mais il est clair par définition des $u_i^*(\mathscr P_\bullet^{(i)})$ et en
vertu de la commutativité de (6.7.10.1) que cet homomorphisme est aussi un
homomorphisme de $(B_1\otimes_A B_2)$-modules; comme le second membre de
(6.7.10.4) est formé de $(B'_1\otimes_A B'_2)$-modules, on déduit
canoniquement de (6.7.10.4) un homomorphisme de
$(B'_1\otimes_A B'_2)$-modules
\[
  {}^{(t)}E
  (X^{(1)},X^{(2)};\mathscr P_\bullet^{(1)},\mathscr P_\bullet^{(2)})
  \otimes_{B_1\otimes_A B_2}(B'_1\otimes_A B'_2)
  \longrightarrow
  {}^{(t)}E
  (X'^{(1)},X'^{(2)};
   u_1^*(\mathscr P_\bullet^{(1)}),u_2^*(\mathscr P_\bullet^{(2)}))
  \tag{6.7.10.5}\label{III.6.7.10.5.fr}
\]
ce qui, compte tenu de (I, 1.6.5) n'est autre que l'homomorphisme cherché
(6.7.10.2) dans le cas particulier considéré.

Il reste à passer au cas général en suivant les recollements de (6.7.1) et
(6.7.2); le second passage est immédiat; en ce qui concerne le premier, on
considère comme dans (6.7.1) des éléments $g_i\in B_i$, et leurs images
$g'_i\in B'_i$, le produit tensoriel $g=g_1\otimes g_2$ dans
$B=B_1\otimes_A B_2$ et son image $g'=g'_1\otimes g'_2$ dans
$B'=B'_1\otimes_A B'_2$, et tout revient à utiliser

